% Формулы метрик и схема данных — Реальт (клиника).
%
% Источник правды для формул — src/schema_realt_metrics_views.sql
% (ClickHouse VIEW) и src/metabase_queries/realt_*.sql (native SQL
% карточек Metabase). Этот файл описывает то же самое человеческим
% языком + формулами, ничего не пересчитывает заново — при изменении
% формулы в VIEW/карточке нужно поправить и здесь (см. последний раздел).
%
% Компиляция: движок XeLaTeX (не pdflatex — нужен fontspec) дважды подряд
% (для оглавления) — xelatex realt.tex && xelatex realt.tex, либо
% latexmk -xelatex realt.tex, либо tectonic realt.tex (автоматически
% нужное число раз). Overleaf: меню Menu -> Compiler -> XeLaTeX.
%
% Шрифты: \mainfont "Times New Roman", \monofont "Courier New" — оба
% протестированы на macOS (полный набор начертаний, кириллица есть в
% обоих). Если на вашей машине/Overleaf такого имени нет — замените
% на "Liberation Serif"/"Liberation Mono" или "TeX Gyre Termes"/"TeX
% Gyre Cursor" (свободные метрически близкие альтернативы, обычно есть
% в любом TeX Live). Проверять кириллицу в MONO-шрифте отдельно — у
% умолчательного (lmmono, Latin Modern Mono) её нет вообще, текст вроде
% "ФОТ ШАА" внутри \texttt печатает пустые прямоугольники вместо букв.
%
% Проверено полной сборкой (tectonic 0.17.0, движок XeTeX) 2026-09-18 —
% 11 страниц, 0 ошибок, 0 отсутствующих глифов, ER-диаграмма и все
% таблицы/формулы проверены визуально (постранично, через PyMuPDF) —
% не только "компилируется без ошибок", но и "выглядит правильно".
%
% 2026-09-19: содержание правилось (переход revenue на department='ШАА'
% с учётом периода, full_revenue, new_client_revenue[_shaa],
% new_client_visits_3m[_avg], новый раздел про "За год" для ratio-метрик;
% отдельным заходом в тот же день — новый раздел 5.4 "Накладные расходы и
% маркетинг, распределённые на визит/клиента") в среде БЕЗ LaTeX-движка
% (tectonic/xelatex/latexmk недоступны ни разу за день) — визуальная
% проверка сборкой НЕ переделана ни для одной из этих правок. Перед тем
% как считать документ актуальным для клиента — пересобрать и визуально
% проверить (см. команду выше), особенно новые формулы, обновлённый
% раздел 4 и новый раздел 5.4 (много новых \texttt{\seqsplit{...}}-идентификаторов
% и вложенных \code{...} — риск переполнения строки/таблицы выше обычного).
% Единственное, что стоит перепроверить при заметной правке текста —
% длинные \texttt{\seqsplit{...}}-идентификаторы (SQL/ClickHouse-имена
% без пробелов) при необходимости переносятся посимвольно без дефиса;
% \code{...} (без seqsplit) используют только фрагменты с пробелами
% внутри (короткие условия вида "amount > 0") — они короче и обычно не
% требуют переноса.

\documentclass[11pt,a4paper]{article}
\usepackage{fontspec}
\setmainfont{Times New Roman}
\setmonofont{Courier New} % умолчательный моноширинный (lmmono) без кириллицы — \texttt{\seqsplit{}} с русскими значениями (например, "ФОТ ШАА") иначе печатает пустые прямоугольники
\usepackage[russian]{babel}
\usepackage[margin=2.2cm]{geometry}
\usepackage{amsmath,amssymb}
\usepackage{xcolor}
\usepackage{booktabs}
\usepackage{longtable}
\usepackage{array}
\usepackage{graphicx}
\usepackage{enumitem}
\usepackage{seqsplit} % позволяет переносить длинные SQL/ClickHouse-идентификаторы без пробелов (например, realt_visits_categorized) внутри строки — иначе они вылезают за поля
\usepackage{tikz}
\usetikzlibrary{positioning, fit, arrows.meta}
\usepackage[colorlinks=true, linkcolor=blue!50!black, urlcolor=blue!50!black, citecolor=blue!50!black]{hyperref}

% Печать SQL/ClickHouse-идентификаторов моноширинным шрифтом без ручного
% экранирования подчёркиваний ("_" в обычном тексте LaTeX — спецсимвол).
\newcommand{\code}[1]{\texttt{\detokenize{#1}}}
\newcolumntype{L}[1]{>{\raggedright\arraybackslash}p{#1}}


\title{Реальт --- формулы метрик и схема данных}
\author{bf-analytics-platform}
\date{\today}

\begin{document}
\maketitle
\tableofcontents
\clearpage

\section{Как пользоваться документом}

Единое место, где видно, \emph{как именно} считается каждая метрика на
дашбордах Metabase проекта «Реальт» (клиника), и как связаны между собой
таблицы-источники. Формулы здесь --- не отдельная версия правды, а
человекочитаемое описание того, что реально лежит в ClickHouse VIEW
(\texttt{\seqsplit{src/ schema\_realt\_metrics\_views.sql}}) и native SQL карточек
(\texttt{\seqsplit{src/ metabase\_queries/ realt\_*.sql}}). Если формула в коде меняется ---
этот документ нужно обновить тем же коммитом (см.
раздел~\ref{sec:maintain}).

Структура: раздел~\ref{sec:terms} --- термины и ключи, которые
используются везде дальше; разделы~\ref{sec:sources}--\ref{sec:other-views}
--- сами метрики, сгруппированные по VIEW/карточке, от которых они
считаются; раздел~\ref{sec:filters} --- параметры (фильтры) дашбордов;
раздел~\ref{sec:er} --- схема соединения таблиц-источников.

\section{Термины и ключи}
\label{sec:terms}

\begin{longtable}{@{}L{3.6cm}L{11cm}@{}}
\toprule
Термин & Определение \\
\midrule
\endfirsthead
\toprule
Термин & Определение \\
\midrule
\endhead
\bottomrule
\endfoot
\textbf{Визит} &
Одна строка \texttt{\seqsplit{klientiks\_operations}} с суммой $>0$ и заполненной датой
начала (\texttt{\seqsplit{visit\_start}}) --- один приём/услуга одного клиента. \\
\textbf{Клиент} &
Идентифицируется по \texttt{\seqsplit{card\_number}} (номер карты в Клиентикс) ---
стабильный ID для когорт/LTV, не по ФИО/телефону. \\
\textbf{Врач / сотрудник} &
Две разные вещи, связанные ключом. В \texttt{\seqsplit{klientiks\_operations.doctor}}
--- текстовое ФИО (как ввели в Клиентикс). В \texttt{\seqsplit{realt\_payroll}} и
\texttt{\seqsplit{realt\_employees}} --- \texttt{\seqsplit{employee\_id}}, короткий код (например,
\texttt{\seqsplit{АрсТБ\_ТД}}). Связка ФИО~$\to$~\texttt{\seqsplit{employee\_id}} --- через
\texttt{\seqsplit{realt\_employees.full\_name}}, сравнение
\texttt{\seqsplit{upperUTF8(trim(...))}} (регистронезависимо, без пробелов по краям).
См.\ раздел~\ref{sec:categorized}. \\
\textbf{Роль / \texttt{\seqsplit{role\_group}}} &
Категория ФОТ сотрудника: ФОТ Психиатры / ФОТ Психологи / ФОТ
Администраторы / ФОТ Управление / ФОТ Маркетинг / \textbf{ФОТ ШАА}.
Начиная с 2026-09-17 определяется по \code{department = 'ШАА'} в
\texttt{\seqsplit{realt\_payroll}} (личность сотрудника), а не по тексту услуги --- см.
раздел~\ref{sec:categorized}. \\
\textbf{ШАА} &
Историческое сокращение по фамилиям Шмилович/Онегина (изначально ---
подразделение этих двух врачей). Сейчас как \texttt{\seqsplit{role\_group}} включает
также их администратора (Ерзаулова Анастасия) --- см.\ таблицу~\ref{tab:role-group}. \\
\textbf{Новый клиент (когорта)} &
Клиент, чей \emph{первый визит за всю историю} пришёлся на рассматриваемый
период. Считается по вычисленному номеру визита
\code{visit_seq = row_number() OVER (PARTITION BY card_number ORDER BY visit_start)}
--- по ВСЕЙ истории клиента, а не по видимым в текущем фильтре визитам.
\textbf{Не путать} с фильтром дашборда «Тип клиента: Первичный/Повторный»
--- см.\ таблицу~\ref{tab:dash-filters}, это про визиты внутри одного
месяца, другая метрика с похожим названием. \\
\end{longtable}

\section{Источники данных}
\label{sec:sources}

Сырые ClickHouse-таблицы, из которых считаются все метрики ниже. Полные
схемы --- \texttt{\seqsplit{src/ schema\_klientiks.sql}} и
\texttt{\seqsplit{src/ schema\_realt\_gsheets.sql}}.

\begin{itemize}[leftmargin=*]
\item \textbf{\texttt{\seqsplit{klientiks\_operations}}} --- построчная выгрузка из
  учётной системы клиники (Клиентикс): один визит = одна строка
  (клиент/врач/услуга/дата/сумма).
\item \textbf{\texttt{\seqsplit{realt\_payroll}}} --- вкладка «Импорт ФОТ» Google-Таблицы:
  начисления по сотрудникам помесячно (роль, тип оплаты, начислено,
  НДФЛ, взносы).
\item \textbf{\texttt{\seqsplit{realt\_employees}}} --- справочник
  \texttt{\seqsplit{employee\_id}}~$\leftrightarrow$~ФИО, 84 строки без дублей.
  \textit{Внимание для поддержки репозитория:} эта таблица и
  \texttt{\seqsplit{realt\_service\_categories}} задеплоены и загружены на прод, но код
  их создания (\texttt{\seqsplit{feat/ realt-employees}}, \texttt{\seqsplit{feat/ realt-service-categories}})
  на момент написания документа \textbf{не смёржен в master} --- см.\
  раздел~\ref{sec:maintain}.
\item \textbf{\texttt{\seqsplit{realt\_service\_categories}}} --- справочник
  \texttt{\seqsplit{service}}~$\to$~категории (тип врача, квалификация и т.д.), ключ ---
  \texttt{\seqsplit{service}} (как в \texttt{\seqsplit{klientiks\_operations.service}}).
  \textbf{Пока не используется} ни в одной VIEW/карточке из этого
  документа --- загружен про запас, для будущего среза «уровень врача».
\item \textbf{\texttt{\seqsplit{realt\_expenses}}} --- вкладка «Остальные расходы»:
  месяц~$\times$~статья, суммы отрицательные (расход).
\item \textbf{\texttt{\seqsplit{realt\_bank\_account}} / \texttt{\seqsplit{realt\_cash}} / \texttt{\seqsplit{realt\_accruals}}}
  --- вкладки «Расчетный счет» / «Наличные» / «Начисления»: плоские
  регистры операций (по кассовому методу и по методу начисления).
  Между собой НЕ джойнятся --- объединяются \code{UNION ALL} по
  одинаковой роли колонок (\texttt{\seqsplit{pl\_article}}, дата, сумма) в
  \texttt{\seqsplit{realt\_pl\_by\_group\_month}} (раздел~\ref{sec:pl}).
\end{itemize}

\section{Базовые фильтры видимости визита}
\label{sec:base-filters}

\textbf{Изменено 2026-09-19} --- раньше этот раздел описывал текстовое
исключение Шмиловича/Онегиной прямо в \texttt{\seqsplit{realt\_metrics\_by\_month}}. Теперь
компания-целиком метрики (раздел~\ref{sec:monthly}) читают визиты \textbf{не
напрямую} из \texttt{\seqsplit{klientiks\_operations}}, а через
\texttt{\seqsplit{realt\_visits\_categorized}} (раздел~\ref{sec:categorized}) --- ту же VIEW,
которой уже пользовалась юнит-экономика по врачу. Базовые фильтры теперь
живут \textbf{там}:

\begin{itemize}[leftmargin=*]
\item \code{amount > 0} --- нулевые/отрицательные суммы не считаются
  выручкой;
\item \texttt{\seqsplit{visit\_start}} не пусто;
\item \code{card_number != ''};
\item \code{positionCaseInsensitiveUTF8(doctor, 'тест') = 0} --- исполнитель
  «тест» исключён;
\item \code{role_group != 'ФОТ ШАА'} --- визиты сотрудников подразделения
  ШАА исключены из выручки/визитов/клиентов компании целиком.
\end{itemize}

\textbf{Исключение ШАА теперь личностное, не текстовое} (до 2026-09-19
было \code{positionCaseInsensitiveUTF8(service, 'шмил'/'онег') = 0} --- по
названию услуги). Раздел~\ref{sec:categorized} описывает, как именно
\texttt{\seqsplit{role\_group}} определяется по \texttt{\seqsplit{employee\_id}} --- \textbf{эта же логика
теперь используется и здесь, и в юнит-экономике по врачу
(раздел~\ref{sec:doctor-econ}), метод уже не «сознательно разный», как
было раньше} (см.\ также раздел~\ref{sec:categorized-diff}).

\textbf{Почему поменяли (найдено при сверке с клиентской отчётностью,
2026-09-19):} у текстового способа было два изъяна --- визит Шмиловича с
обычным названием услуги не исключался, а чужой визит с его именем в
составном названии услуги, наоборот, ложно исключался (расхождение на
реальных данных --- около 2.46\,млн~руб. / 313 визитов). Но личностный
способ в его \emph{первой} версии (2026-09-17, \code{DISTINCT employee_id
FROM realt_payroll WHERE department = 'ШАА'}, без учёта периода) оказался
\textbf{тоже неверным} для сотрудников, у которых \code{department}
менялся во времени: Онегина Елена Юрьевна была в отделе «3 этаж»
(роль «ФОТ Психиатры») по март 2026 включительно и переведена в ШАА
только с апреля 2026 --- старый запрос относил к ШАА \emph{весь} её стаж
задним числом, занижая выручку компании за январь--март 2026 ровно на её
выручку за эти месяцы (561\,000 / 510\,500 / 347\,000~руб. --- сверено с
клиентской таблицей до рубля). Исправлено на привязку \code{department} к
\textbf{конкретному месяцу} периода ФОТ (\code{JOIN} по
\texttt{\seqsplit{employee\_id}} + месяц, не просто \texttt{\seqsplit{employee\_id}}) --- см.\
\texttt{\seqsplit{payroll\_period}} в \texttt{\seqsplit{realt\_visits\_categorized}}
(раздел~\ref{sec:categorized}).

\section{Месячные метрики компании (\texttt{realt\_metrics\_by\_month})}
\label{sec:monthly}

VIEW: одна строка = один месяц (\texttt{\seqsplit{month}}, заякорен на 12:00 --- иначе
Report Timezone Metabase сдвигает 1-е число на конец предыдущего месяца).
Строится \code{FULL OUTER JOIN} нескольких агрегатов по месяцу --- визиты
без ШАА, ФОТ, полная выручка (все визиты), визиты ШАА, скользящая
когорта (см.\ ниже) --- сознательно не \code{LEFT JOIN} с визитов,
иначе месяц, где ФОТ уже загружен, а визитов в Клиентикс ещё нет
(например, текущий месяц), молча пропадает из VIEW целиком (это был баг,
исправленный 2026-09-17).

\subsection{Выручка, визиты, клиенты}

Источник --- \texttt{\seqsplit{realt\_visits\_categorized}} (раздел~\ref{sec:categorized})
\textbf{с фильтром} \code{role_group != 'ФОТ ШАА'} (см.\
раздел~\ref{sec:base-filters}). \texttt{\seqsplit{visit\_seq}} здесь --- ПЕРЕСЧИТАННЫЙ
номер визита клиента \emph{после} исключения ШАА (\code{row_number()} по
уже отфильтрованному набору), а не общий \texttt{\seqsplit{visit\_seq}} из
\texttt{\seqsplit{realt\_visits\_categorized}} --- так «новый клиент» по-прежнему значит
«первый визит к клинике без учёта ШАА», а не «первый визит вообще
куда-либо, включая ШАА»:

\[
  \texttt{\seqsplit{revenue}} = \sum \texttt{\seqsplit{amount}}
  \qquad
  \texttt{\seqsplit{visits}} = \text{count}(*)
  \qquad
  \texttt{\seqsplit{clients}} = \text{uniqExact}(\texttt{\seqsplit{card\_number}})
\]
\[
  \texttt{\seqsplit{new\_clients}} = \text{uniqExactIf}(\texttt{\seqsplit{card\_number}},\; \texttt{\seqsplit{visit\_seq}} = 1)
  \qquad
  \texttt{\seqsplit{avg\_check}} = \dfrac{\texttt{\seqsplit{revenue}}}{\texttt{\seqsplit{visits}}}
\]

По визитам месяца, прошедшим базовые фильтры (раздел~\ref{sec:base-filters}).
\texttt{\seqsplit{new\_clients}} --- см.\ определение «Новый клиент» в
таблице терминов, раздел~\ref{sec:terms}. \texttt{\seqsplit{avg\_check}} ---
\textbf{ratio-колонка}, см.\ раздел~\ref{sec:ratio-yearly} про то, почему
её нельзя суммировать по месяцам для годового итога.

\subsection{Полная выручка, ШАА и выручка с первого визита (добавлено 2026-09-19)}
\label{sec:new-revenue-cols}

\[
  \texttt{\seqsplit{full\_revenue}} = \sum \texttt{\seqsplit{amount}} \text{ по ВСЕМ визитам месяца (}\texttt{\seqsplit{realt\_visits\_categorized}}\text{, БЕЗ фильтра по } \texttt{\seqsplit{role\_group}}\text{)}
\]

Единственная колонка этой VIEW, где ШАА \textbf{не} исключается ---
«полная выручка компании», по просьбе клиента, для сравнения с
\texttt{\seqsplit{revenue}} (без ШАА).

\[
  \texttt{\seqsplit{new\_client\_revenue}} = \sum \texttt{\seqsplit{amount}} \text{ где } \texttt{\seqsplit{visit\_seq}} = 1
  \qquad\text{(тот же }\texttt{\seqsplit{visit\_seq}}\text{, что и у }\texttt{\seqsplit{new\_clients}}\text{, без ШАА)}
\]

Выручка \textbf{именно} за первый визит новых клиентов месяца (не вся их
последующая выручка). До 2026-09-19 называлась «Выручка с первых
визитов» --- переименована в Metabase в «Выручка с 1 визита», формула не
менялась.

\[
  \texttt{\seqsplit{new\_clients\_shaa}} = \text{uniqExactIf}(\texttt{\seqsplit{card\_number}},\; \texttt{\seqsplit{visit\_seq}}_{\text{ШАА}} = 1)
  \qquad
  \texttt{\seqsplit{new\_client\_revenue\_shaa}} = \sum \texttt{\seqsplit{amount}} \text{ где } \texttt{\seqsplit{visit\_seq}}_{\text{ШАА}} = 1
\]

Те же формулы, что \texttt{\seqsplit{new\_clients}}/\allowbreak\texttt{\seqsplit{new\_client\_revenue}}, но
для визитов ШАА (\code{role_group = 'ФОТ ШАА'}) и с СОБСТВЕННОЙ
последовательностью \texttt{\seqsplit{visit\_seq}}$_{\text{ШАА}}$, посчитанной только по
визитам ШАА (\texttt{\seqsplit{shaa\_visits\_seq}} в VIEW). \textbf{Независимая когорта} от
\texttt{\seqsplit{new\_clients}}: клиент мог годами наблюдаться в клинике и прийти к ШАА
впервые (войдёт в \texttt{\seqsplit{new\_clients\_shaa}}, но не в \texttt{\seqsplit{new\_clients}}), и
наоборот (новый в клинике, но не был у ШАА).

\subsection{Визиты новых клиентов за скользящий квартал (добавлено 2026-09-19)}
\label{sec:rolling-visits}

Когорта --- «Новые клиенты» месяца $M$ (раздел выше, БЕЗ ШАА). Для
каждого такого клиента считаем ВСЕ его визиты (тоже без ШАА), которые
попали в 3-месячное \textbf{скользящее} окно $[M \dots M{+}2]$ включительно
(месяц привлечения и два следующих) --- \textbf{не} выравнивание по
календарным кварталам: клиент, пришедший в феврале, считается за
февраль--март--апрель, а не за Q1/Q2.

\[
  \texttt{\seqsplit{new\_client\_visits\_3m}} = \sum_{\text{клиенты, привлечённые в } M} \big(\text{их визиты в } [M \dots M{+}2]\big)
\]
\[
  \texttt{\seqsplit{new\_client\_visits\_3m\_avg}} = \frac{\texttt{\seqsplit{new\_client\_visits\_3m}}}{\texttt{\seqsplit{new\_clients}}}
\]

\textbf{Важно про свежие месяцы:} для последних 1--2 загруженных месяцев
окно уходит за пределы того, что уже есть в базе (например, для когорты
июля 2026 нужны визиты августа и сентября, которых ещё нет) --- значение
занижено \emph{технически}, это не падение активности новых клиентов.
\texttt{\seqsplit{new\_client\_visits\_3m\_avg}} --- \textbf{ratio-колонка}, см.\
раздел~\ref{sec:ratio-yearly}.

\subsection{ФОТ по ролям и типу оплаты}

Общая формула для \textbf{любой} строки таблицы~\ref{tab:fot-monthly}:
сумма \texttt{\seqsplit{accrued\_total}} (колонка «Начислено ИТОГО» из
\texttt{\seqsplit{realt\_payroll}}, \textbf{не} «К оплате» --- согласовано с клиентом
2026-09-14) по строкам месяца, отфильтрованным по роли и/или типу оплаты:

\[
  \text{ФОТ}_{\text{роль},\,\text{тип}} = \sum \texttt{\seqsplit{accrued\_total}}
  \quad\text{где}\quad \texttt{\seqsplit{role}} = \text{роль} \;\wedge\; \texttt{\seqsplit{pay\_type}} = \text{тип}
\]

\begin{longtable}{@{}L{5.2cm}L{5.5cm}L{4cm}@{}}
\caption{Колонки ФОТ в \texttt{realt\_metrics\_by\_month} (все --- частные
случаи формулы выше; «---» = без фильтра по этому измерению)}
\label{tab:fot-monthly}\\
\toprule
Колонка & Роль (\texttt{\seqsplit{role}}) & Тип оплаты (\texttt{\seqsplit{pay\_type}}) \\
\midrule
\endfirsthead
\toprule
Колонка & Роль (\texttt{\seqsplit{role}}) & Тип оплаты (\texttt{\seqsplit{pay\_type}}) \\
\midrule
\endhead
\bottomrule
\endfoot
\texttt{\seqsplit{fot\_total}} & --- (все роли) & --- \\
\texttt{\seqsplit{fot\_total\_pct}} & --- & Проценты \\
\texttt{\seqsplit{fot\_total\_oklad}} & --- & Оклад, Бонус \\
\texttt{\seqsplit{fot\_psychiatrists}} & ФОТ Психиатры & --- \\
\texttt{\seqsplit{fot\_psychiatrists\_pct}} & ФОТ Психиатры & Проценты \\
\texttt{\seqsplit{fot\_psychiatrists\_oklad}} & ФОТ Психиатры & Оклад, Бонус \\
\texttt{\seqsplit{fot\_psychologists}} & ФОТ Психологи & --- \\
\texttt{\seqsplit{fot\_psychologists\_pct}} & ФОТ Психологи & Проценты \\
\texttt{\seqsplit{fot\_psychologists\_oklad}} & ФОТ Психологи & Оклад, Бонус \\
\texttt{\seqsplit{fot\_administrators}} & ФОТ Администраторы & --- \\
\texttt{\seqsplit{fot\_administrators\_pct}} & ФОТ Администраторы & Проценты \\
\texttt{\seqsplit{fot\_administrators\_oklad}} & ФОТ Администраторы & Оклад, Бонус \\
\texttt{\seqsplit{fot\_management}} & ФОТ Управление & --- \\
\texttt{\seqsplit{fot\_management\_pct}} & ФОТ Управление & Проценты \\
\texttt{\seqsplit{fot\_management\_oklad}} & ФОТ Управление & Оклад, Бонус \\
\texttt{\seqsplit{fot\_marketing}} & ФОТ Маркетинг & --- \\
\texttt{\seqsplit{fot\_marketing\_pct}} & ФОТ Маркетинг & Проценты \\
\texttt{\seqsplit{fot\_marketing\_oklad}} & ФОТ Маркетинг & Оклад, Бонус \\
\texttt{\seqsplit{fot\_shmilovich}} & \code{department = 'ШАА'} \newline (не роль, а подразделение!) & --- \\
\texttt{\seqsplit{fot\_shmilovich\_pct}} & \code{department = 'ШАА'} & Проценты \\
\texttt{\seqsplit{fot\_shmilovich\_oklad}} & \code{department = 'ШАА'} & Оклад, Бонус \\
\end{longtable}

\textbf{\texttt{\seqsplit{fot\_shmilovich}} --- не то же самое, что раньше.} До
2026-09-17 ключ был \code{role = 'ФОТ Шмилович'} (33 строки, $\approx$22.97\,млн~руб.,
строго сам Шмилович). Сейчас --- \code{department = 'ШАА'} (49 строк,
$\approx$24.27\,млн~руб.): шире, включает ещё администратора/психиатра,
числящихся в подразделении ШАА. Это меняет ранее показанные клиенту
цифры «Доля ФОТ в выручке» / «Прибыль после ФОТ» (согласованные
2026-09-14 по старому ключу) --- упомянуть при следующей сверке.

\subsection{Налоги и доля ФОТ в выручке}

\[
  \texttt{\seqsplit{fot\_taxes}} = \sum \texttt{\seqsplit{ndfl}} + \sum \texttt{\seqsplit{contributions}}
\]

Значения в источнике обычно \textbf{отрицательные} (это удержания, не
начисления). \texttt{\seqsplit{fot\_taxes\_pct}} / \texttt{\seqsplit{fot\_taxes\_oklad}} --- то же с
фильтром по \texttt{\seqsplit{pay\_type}}, как в таблице~\ref{tab:fot-monthly}.

\[
  \texttt{\seqsplit{fot\_revenue\_share}} = \frac{\texttt{\seqsplit{fot\_total}} - \texttt{\seqsplit{fot\_shmilovich}}}{\texttt{\seqsplit{revenue}}} \times 100
\]

ФОТ подразделения ШАА исключён из \textbf{числителя}, потому что выручка
Шмиловича/Онегиной \textbf{не входит} в \texttt{\seqsplit{revenue}} (текстовое
исключение, раздел~\ref{sec:base-filters}) --- иначе доля была бы
искусственно завышена. \texttt{\seqsplit{fot\_total}} при этом остаётся полным (не
вычитается из него ШАА нигде, кроме этой формулы).

\subsection{Прибыль после ФОТ (уровень Metabase Model, не VIEW)}

\[
  \text{Прибыль после ФОТ} = \texttt{\seqsplit{revenue}} - \texttt{\seqsplit{fot\_total}} + \texttt{\seqsplit{fot\_taxes}}
\]

\textbf{Это не колонка \texttt{\seqsplit{realt\_metrics\_by\_month}}} --- формула
вычисляется в самой Metabase Model («Модель - Реальт метрики по
месяцам», id~98, \texttt{\seqsplit{src/ metabase\_queries/ realt\_metrics\_model.sql}}) поверх
готовых \texttt{\seqsplit{revenue}}/\allowbreak\texttt{\seqsplit{fot\_total}}/\allowbreak\texttt{\seqsplit{fot\_taxes}} из VIEW. Так как
\texttt{\seqsplit{fot\_taxes}} обычно $\le 0$ (удержания), прибавление, а не вычитание
--- корректно, знак уже заложен в данных.

\subsection{«За год» для ratio-метрик --- считается НЕ в VIEW (добавлено 2026-09-19)}
\label{sec:ratio-yearly}

\texttt{\seqsplit{realt\_metrics\_by\_month}} --- помесячная VIEW, у неё вообще нет
годового грейна: колонки вроде \texttt{\seqsplit{avg\_check}} или
\texttt{\seqsplit{new\_client\_visits\_3m\_avg}} --- это отношения (\emph{ratio}), посчитанные
\textbf{для каждого месяца отдельно}. Годовой итог такой колонки --- НЕ
сумма и не среднее из 12 месячных значений:

\[
  \underbrace{\sum_{m=1}^{12} \texttt{\seqsplit{avg\_check}}_m}_{\text{неверно --- сумма 12 разных средних}}
  \quad\ne\quad
  \underbrace{\frac{\sum_{m=1}^{12} \texttt{\seqsplit{revenue}}_m}{\sum_{m=1}^{12} \texttt{\seqsplit{visits}}_m}}_{\text{верно --- средний чек за весь год}}
\]

На реальных данных 2026 года левая часть --- около 66\,812 (сумма семи
доступных месячных средних), правая --- 9\,540 (реальный средний чек за
период) --- разница на порядок, не округление.

\textbf{Где действительно считается верная годовая формула:} только в
дашборд-карточке «Таблица - Реальт - Ежемесячные метрики» (id~186,
\texttt{\seqsplit{src/ metabase\_queries/ realt\_monthly\_metrics.sql}}) --- она вручную
строит столбец «За год» отдельной формулой для каждой ratio-строки
(\code{sum(Выручка)/sum(Визиты)} для среднего чека,
\code{sum(new_client_visits_3m)/sum(new_clients)} для визитов на нового
клиента), а не переиспользует \texttt{\seqsplit{sumIf}} по месяцам, как для
аддитивных метрик (Выручка, Визиты, Новые клиенты и т.д.). \textbf{В самой
ClickHouse VIEW такой годовой агрегации нет вообще} --- если понадобится
годовое значение где-то ещё (например, боту, который не пойдёт в эту
конкретную Metabase-карточку, а прочитает VIEW напрямую), формулу нужно
повторить \emph{осознанно}, а не через \code{avg()}/\code{sum()} по
\texttt{\seqsplit{month}}.

Колонки-ratio в этой VIEW на 2026-09-19: \texttt{\seqsplit{avg\_check}},
\texttt{\seqsplit{new\_client\_visits\_3m\_avg}}, \texttt{\seqsplit{fot\_revenue\_share}} (раздел выше ---
хоть она и не выведена в карточку 186, тот же принцип применим, если её
туда добавят). У каждой такой колонки \code{COMMENT COLUMN} в
\texttt{\seqsplit{schema\_realt\_metrics\_views.sql}} явно предупреждает об этом же
(см.\ \code{architecture-standarts.md} $\to$ «Семантический слой для
AI-бота» --- то же требование для бота, читающего \texttt{\seqsplit{system.columns}}
в обход этого документа).

\textbf{Это же причина, по которой карточка 186 --- Таблица, а не
Визуал/MBQL-пивот}, и осознанное исключение из правила «на дашборд ---
только Визуалы» (\code{architecture-standarts.md}): MBQL-пивот считает
столбец «Итого» ТОЙ ЖЕ агрегацией, что и ячейки (\code{sum}), у него нет
способа посчитать для одной строки-метрики \code{sum}, а для другой ---
\code{sum/sum}.

\section{Категоризация визита и ФОТ по врачу}
\label{sec:categorized}

Две VIEW-надстройки над \texttt{\seqsplit{klientiks\_operations}} и \texttt{\seqsplit{realt\_payroll}}
--- добавляют \texttt{\seqsplit{employee\_id}} и \texttt{\seqsplit{role\_group}} \textbf{на уровне
конкретного визита/строки зарплаты}, а не только помесячного агрегата.
Именно отсюда берётся юнит-экономика по конкретному врачу
(раздел~\ref{sec:doctor-econ}).

\subsection{Ключ сопоставления врач \texorpdfstring{$\leftrightarrow$}{<->} сотрудник}

\begin{quote}
\texttt{\seqsplit{upperUTF8(trim(klientiks\_operations.doctor))}} = \texttt{\seqsplit{upperUTF8(trim(realt\_employees.full\_name))}}
\end{quote}

Регистронезависимое сравнение по ФИО (\texttt{\seqsplit{upperUTF8}}, не \texttt{\seqsplit{upper}}
--- обычный \texttt{\seqsplit{upper()}} не работает с кириллицей). На момент внедрения
(2026-09-17) покрытие --- 100\,\% с учётом регистра; 6 записей остаются с
разным регистром ФИО в справочнике сотрудников (не критично).

\subsection{\texttt{role\_group} --- приоритет правил (\texttt{multiIf})}
\label{sec:role-group-logic}

\begin{longtable}{@{}cL{4.8cm}L{7.3cm}@{}}
\caption{Порядок проверки условий в \texttt{realt\_visits\_categorized}
(первое совпавшее условие побеждает)}
\label{tab:role-group}\\
\toprule
\# & Условие & Результат (\texttt{\seqsplit{role\_group}}) \\
\midrule
\endfirsthead
\toprule
\# & Условие & Результат (\texttt{\seqsplit{role\_group}}) \\
\midrule
\endhead
\bottomrule
\endfoot
1 & Врач найден в \texttt{\seqsplit{realt\_employees}} \textbf{и} за МЕСЯЦ ЭТОГО ВИЗИТА
    (не когда-либо) у его \texttt{\seqsplit{employee\_id}} есть хотя бы одна строка
    \texttt{\seqsplit{realt\_payroll}} с \code{department = 'ШАА'} (\code{max(department =
    'ШАА')} по строкам месяца --- см.\ оговорку про Ерзаулову ниже)
  & \textbf{ФОТ ШАА} \\
2 & Врач найден в \texttt{\seqsplit{realt\_employees}} \textbf{и} для его
    \texttt{\seqsplit{employee\_id}} есть строка \texttt{\seqsplit{realt\_payroll}} ЗА ТОТ ЖЕ МЕСЯЦ
    (\code{argMax(role, period)} внутри месяца, если строк несколько)
  & эта роль (ФОТ Психиатры / Психологи / Администраторы / Управление / Маркетинг) \\
3 & Врач найден в \texttt{\seqsplit{realt\_employees}}, но \texttt{\seqsplit{employee\_id}} \textbf{ни разу}
    не встречается в \texttt{\seqsplit{realt\_payroll}} ЗА МЕСЯЦ ЭТОГО ВИЗИТА (в т.ч. если
    вообще нет ни одной строки ФОТ за этот месяц --- напр. текущий месяц
    до того, как ФОТ загружен)
  & \code{Без ФОТ-кода} (старый/бывший сотрудник вне системы ФОТ, либо ФОТ
    месяца ещё не загружен --- ожидаемо, не баг) \\
4 & ФИО из Клиентикс вообще не нашло пару в \texttt{\seqsplit{realt\_employees}}
  & \code{Не в справочнике сотрудников} \\
\end{longtable}

\textbf{Изменено 2026-09-19 --- привязка к месяцу, а не «когда-либо».}
Первая версия (2026-09-17) брала \code{DISTINCT employee_id FROM
realt_payroll WHERE department = 'ШАА'} \textbf{без учёта периода} --- то
есть если сотрудник хоть раз был в ШАА, к ШАА относился весь его стаж,
включая визиты ДО перевода. На реальных данных это ошибочно отнесло к ШАА
январь--март 2026 Онегиной Е.Ю. (тогда --- «3 этаж», роль «ФОТ
Психиатры»; в ШАА --- только с апреля 2026), см.\
раздел~\ref{sec:base-filters}. Один сотрудник может иметь \textbf{несколько
строк ФОТ за один месяц с разным \texttt{\seqsplit{department}}} (например,
администратор Ерзаулова --- частичная занятость и на этаже, и в ШАА) ---
месяц считается «ШАА», если ШАА встретилась хотя бы в одной его строке
ФОТ этого месяца.

На реальных данных \code{role_group = 'ФОТ ШАА'} сейчас --- Шмилович
Андрей Аркадьевич (ШАА весь период), Онегина Елена Юрьевна (ШАА только с
апреля 2026) и их администратор Ерзаулова Анастасия (частичная занятость
в ШАА каждый месяц). \texttt{\seqsplit{realt\_payroll\_categorized}} использует ту же
логику, но ещё проще --- \texttt{\seqsplit{department}} уже есть в каждой строке
\texttt{\seqsplit{realt\_payroll}} за её собственный период, доп.\ \texttt{\seqsplit{JOIN}} не нужен:
\[
  \texttt{\seqsplit{role\_group}} = \begin{cases}
    \text{ФОТ ШАА}, & \texttt{\seqsplit{department}} = \text{'ШАА'}\ (\text{этой же строки})\\
    \texttt{\seqsplit{role}} \text{ (как есть)}, & \text{иначе}
  \end{cases}
\]

\subsection{Чем это отличается от раздела~\ref{sec:base-filters}}
\label{sec:categorized-diff}

\textbf{Изменено 2026-09-19 --- больше не отличается.} До этой даты
компания-целиком метрики (\texttt{\seqsplit{realt\_metrics\_by\_month}}) исключали
Шмиловича/Онегину \textbf{по тексту названия услуги}, а здесь --- \textbf{по
личности врача} (через \texttt{\seqsplit{employee\_id}}). С 2026-09-19
\texttt{\seqsplit{realt\_metrics\_by\_month}} тоже читает визиты через
\texttt{\seqsplit{realt\_visits\_categorized}} и исключает по \code{role_group != 'ФОТ
ШАА'} --- \textbf{тот же личностный способ}, что и здесь, единый ключ для
выручки компании и для юнит-экономики по врачу. Требование --- чтобы врач
был найден в \texttt{\seqsplit{realt\_employees}}, иначе визит попадёт в категорию
«Без ФОТ-кода» или «Не в справочнике сотрудников» (таблица~\ref{tab:role-group})
--- эти визиты \textbf{не исключаются} из выручки компании (только ШАА
исключается), просто не участвуют в разбивке по ролям.

\section{Юнит-экономика по врачу и по роли}
\label{sec:doctor-econ}

\texttt{\seqsplit{realt\_doctor\_month}} (месяц~$\times$~врач) и \texttt{\seqsplit{realt\_role\_month}}
(месяц~$\times$~роль) --- агрегаты поверх
\texttt{\seqsplit{realt\_visits\_categorized}}/\allowbreak\texttt{\seqsplit{realt\_payroll\_categorized}}
(раздел~\ref{sec:categorized}). Формулы идентичны для обеих VIEW,
отличается только уровень группировки (\texttt{\seqsplit{employee\_id}} vs
\texttt{\seqsplit{role\_group}}):

\begin{quote}
\texttt{\seqsplit{revenue}}, \texttt{\seqsplit{visits}}, \texttt{\seqsplit{clients}}, \texttt{\seqsplit{new\_clients}}
--- те же формулы, что в разделе~\ref{sec:monthly}, но по визитам ЭТОГО врача/роли.

\texttt{\seqsplit{fot}}, \texttt{\seqsplit{fot\_pct}}, \texttt{\seqsplit{fot\_oklad}}, \texttt{\seqsplit{fot\_taxes}}
--- те же формулы, что в таблице~\ref{tab:fot-monthly}, по зарплате ЭТОГО врача/роли.
\end{quote}
\[
  \texttt{\seqsplit{avg\_check}} = \frac{\texttt{\seqsplit{revenue}}}{\texttt{\seqsplit{visits}}}
  \qquad
  \texttt{\seqsplit{fot\_per\_visit}} = \frac{\texttt{\seqsplit{fot}}}{\texttt{\seqsplit{visits}}}
  \qquad
  \texttt{\seqsplit{fot\_per\_client}} = \frac{\texttt{\seqsplit{fot}}}{\texttt{\seqsplit{clients}}}
\]
\[
  \texttt{\seqsplit{profit\_after\_fot}} = \texttt{\seqsplit{revenue}} - \texttt{\seqsplit{fot}} + \texttt{\seqsplit{fot\_taxes}}
\]

\textbf{Ключевое отличие от того, что было раньше} (до 2026-09-17):
«ФОТ на визит» считался как (ФОТ \emph{всей роли} за месяц) / (визиты
\emph{всей клиники} за месяц) --- усреднение по больнице, выбрать
конкретного врача было нельзя. Теперь знаменатель --- визиты именно
ЭТОГО врача (\texttt{\seqsplit{realt\_doctor\_month}}) или визиты, \emph{принятые
именно этой ролью} (\texttt{\seqsplit{realt\_role\_month}}), а не визиты всей клиники.
Так же для \texttt{\seqsplit{fot\_per\_client}} в \texttt{\seqsplit{realt\_role\_month}}:
\textbf{внимание} --- если клиент в одном месяце был и у психиатра, и у
психолога, он войдёт в \texttt{\seqsplit{clients}} обеих ролей --- это не то же
самое, что «уникальные клиенты клиники».

Дашборд-карточки «Помесячная юнитка (по врачам, корректная)» (id~183,
\texttt{\seqsplit{realt\_unitka\_by\_doctor.sql}}) и «Врачи (пофамильно)» (id~184,
\texttt{\seqsplit{realt\_doctors\_table.sql}}) --- те же формулы поверх
\texttt{\seqsplit{realt\_visits\_categorized}}/\allowbreak\texttt{\seqsplit{realt\_payroll\_categorized}}
напрямую (не через готовые \texttt{\seqsplit{realt\_doctor\_month}}/\allowbreak\texttt{\seqsplit{realt\_role\_month}}),
с параметрами дашборда (раздел~\ref{sec:filters}) --- значения при
фильтрах по умолчанию совпадают с \texttt{\seqsplit{realt\_doctor\_month}}/\allowbreak\texttt{\seqsplit{realt\_role\_month}}.
«ФОТ на юнит» там --- \texttt{\seqsplit{fot\_per\_visit}} или \texttt{\seqsplit{fot\_per\_client}} в
зависимости от параметра «Показатели на» (переключатель Визит/Клиент).

\subsection{Накладные расходы и маркетинг, распределённые на визит/клиента (добавлено 2026-09-19, переработано в тот же вечер по этажам)}
\label{sec:overhead}

До этого раздела вся «юнит-экономика по врачу» — это \textbf{прямые}
расходы: ФОТ собственной роли врача, привязанный к нему через
\texttt{\seqsplit{employee\_id}}. Но у клиники есть расходы, которые физически
не привязаны ни к одному конкретному врачу — работа администраторов,
управление, аренда/коммуналка, реклама. Раздел описывает НОВЫЕ
строки/колонки (карточка~183: строки 23--26; карточка~184: колонки
«Накладные на юнит»/«Маркетинг на юнит»/«Прибыль после всех затрат») —
существующие строки 12--17/колонка «ФОТ на юнит» карточек 183/184
\textbf{не менялись}, это разные метрики.

\textbf{История одного вечера (важно для чтения формул ниже):} первая
версия (один общеклинический пул, поровну на все визиты клиники) не
сошлась со реальным P\&L владельца при построчной сверке — нашли два
самостоятельных источника ошибки, независимых от методики распределения
как таковой:
\begin{enumerate}[leftmargin=*]
\item \textbf{Контаминация} — «АПДШ» (отдельное направление/юрлицо, не
  клиника) и «Шмилович личное» утекали в статьи P\&L-whitelist наравне с
  клиникой; «Начисления» (\texttt{\seqsplit{realt\_accruals}}) дублировали часть сумм
  «Расчётного счёта» (Р/с) — один платёж считался и целиком в Р/с, и ещё
  раз размазанным по месяцам в Начислениях. Обе проблемы починены в
  самой VIEW \texttt{\seqsplit{realt\_pl\_by\_group\_month}} (раздел~\ref{sec:pl}), а не
  здесь — правки касаются \textbf{любого} потребителя этой VIEW, включая
  старую карточку «П\&Л по группам» (id~179).
\item \textbf{Методика} — владелец показал реальную структуру P\&L: она
  построена \textbf{по этажам} (2 этаж/3 этаж/ШАА), а не по роли
  визита. У каждого этажа — своя выручка, свой ФОТ врачей, СВОИ
  администраторы, своя аренда. Плоский общеклинический пул смешивал
  расходы одного этажа с визитами другого — переписано ниже.
\end{enumerate}

\textbf{Принцип: ставка считается по клинике целиком, атрибуция — по
видимому фильтру.} Числитель (расход) и знаменатель (визиты/выручка
клиники) в разделах~\ref{sec:overhead-fixed}--\ref{sec:overhead-mkt}
ниже считаются \textbf{без фильтра} \texttt{\seqsplit{doctor\_name}}/\allowbreak
\texttt{\seqsplit{doctor\_type}} (только год и тумблер \texttt{\seqsplit{include\_shaa}}) — ставка не
«скачет» от того, какой врач сейчас выбран. Дальше ставка присваивается
\textbf{каждому} визиту месяца по ЕГО этажу (накладные) и ЕГО роли
(маркетинг), и уже \textbf{сумма по видимым в текущем фильтре визитам}
делится на \texttt{\seqsplit{denom}} (визиты/клиенты этого среза) — так число у
конкретного врача отражает его долю клиники, а не обнуляется.

\subsubsection{Аренда и ФОТ Администраторов — по этажу}
\label{sec:overhead-fixed}

Ключ — новая колонка \texttt{\seqsplit{visit\_floor}} в \texttt{\seqsplit{realt\_visits\_categorized}}
(раздел~\ref{sec:categorized}): \texttt{\seqsplit{department}} сотрудника В МЕСЯЦЕ
ВИЗИТА (2 этаж/3 этаж/ШАА) — вычисляется той же связкой
\texttt{\seqsplit{employee\_id}}+месяц, что и \texttt{\seqsplit{role\_group}}, но это
\textbf{независимая} ось: на 2 и 3 этаже работают И психиатры, И
психологи вперемешку, \texttt{\seqsplit{visit\_floor}} — НЕ то же самое, что «Тип
врача». Проверено на реальных данных 2026-09-19: \texttt{\seqsplit{department}} в
\texttt{\seqsplit{realt\_payroll}} принимает значения «2 этаж»/«3 этаж»/«ШАА» (у
психиатров, психологов И администраторов есть свои значения на каждом
из трёх) и отдельно «УК» (управляющая компания — ТОЛЬКО роли «ФОТ
Управление»/«ФОТ Маркетинг», ни разу не на этаже).

\begin{itemize}[leftmargin=*]
\item \textbf{Аренда+коммуналка+санпэдрежим СВОЕГО этажа} (статьи
  «... - 2 этаж»/«... - 3 этаж»/«... - ШАА» из
  \texttt{\seqsplit{realt\_pl\_by\_group\_month}}) делится ТОЛЬКО на визиты этого же
  этажа — не на всю клинику.
\item \textbf{ФОТ Администраторов} — читается напрямую из
  \texttt{\seqsplit{realt\_payroll}} (не через \texttt{\seqsplit{realt\_payroll\_categorized}} — та VIEW
  подменяет \texttt{\seqsplit{role}} на «ФОТ ШАА» для строк с
  \texttt{\seqsplit{department}}=«ШАА», а здесь нужна ИСХОДНАЯ роль «ФОТ
  Администраторы» именно у ШАА-администратора тоже), фильтруется по
  \texttt{\seqsplit{department}} этажа, делится на визиты только этого этажа.
\end{itemize}

\[
  \texttt{\seqsplit{rent\_rate}}_{\text{этаж}} = \frac{\text{Аренда}_{\text{этаж}}}{\texttt{\seqsplit{visits}}_{\text{этаж}}}
  \qquad
  \texttt{\seqsplit{adm\_rate}}_{\text{этаж}} = \frac{-\text{ФОТ Админы}_{\text{этаж}} + \text{их налоги}_{\text{этаж}}}{\texttt{\seqsplit{visits}}_{\text{этаж}}}
\]

\subsubsection{Управление + «мелкая» административка — только 2 и 3 этаж, по выручке}
\label{sec:overhead-upr}

ФОТ Управления + 8 «мелких» статей группы «Административные» (Сервисы и
подписки, Интернет/телефония/почта, Ремонт, Канцелярия, Аутсорс,
Мероприятия, Банковские услуги, Прочие админ расходы) + «Проценты по
кредитам»/«Амортизация» — по решению владельца 2026-09-19 это ОДИН пул,
который делится \textbf{только между 2 и 3 этажом}, пропорционально их
выручке (\textbf{ШАА не участвует вообще} — у него, со слов владельца,
свои отдельные расходы этого рода, не отражённые в этой модели).
«Налог на прибыль / УСН» \textbf{временно исключён из накладных совсем}
— владелец посчитает его отдельно, фиксированным процентом от выручки
или маржи, формула ещё не согласована на 2026-09-19. «Прочие доходы»
исключены как обычно (доход, не расход).

\[
  \text{Пул}_{\text{упр}} = -\text{ФОТ Управление} + \text{его налоги} + \text{Админ-мелочь} + \text{Кредиты}+\text{Амортизация}
\]
\[
  \texttt{\seqsplit{upr\_rate}}_{\text{2эт}} = \frac{\text{Пул}_{\text{упр}} \times \dfrac{\text{revenue}_{\text{2эт}}}{\text{revenue}_{\text{2эт}}+\text{revenue}_{\text{3эт}}}}{\texttt{\seqsplit{visits\_2et}}}
  \qquad
  \texttt{\seqsplit{upr\_rate}}_{\text{3эт}} = \frac{\text{Пул}_{\text{упр}} \times \dfrac{\text{revenue}_{\text{3эт}}}{\text{revenue}_{\text{2эт}}+\text{revenue}_{\text{3эт}}}}{\texttt{\seqsplit{visits\_3et}}}
\]

Итоговая ставка накладных на визит этажа $X\in\{$2эт,3эт$\}$ —
$\texttt{\seqsplit{rent\_rate}}_X + \texttt{\seqsplit{adm\_rate}}_X + \texttt{\seqsplit{upr\_rate}}_X$; для ШАА — только
$\texttt{\seqsplit{rent\_rate}}_{\text{ШАА}} + \texttt{\seqsplit{adm\_rate}}_{\text{ШАА}}$ (без доли Управления).
Знак везде — как в «Прибыль после ФОТ» (раздел~\ref{sec:monthly}): ФОТ
негируется, налоги/P\&L-статьи складываются как есть.

\textbf{Известный открытый вопрос (2026-09-19, не решён):} выручка
\texttt{\seqsplit{visit\_floor}}=«2 этаж»/«3 этаж» по данным (\texttt{\seqsplit{department}} из ФОТ)
расходится с таблицей владельца на $\approx$266\,тыс~₽/мес (3,3\,\% от
оборота, зеркально — 2 этаж занижен, 3 этаж завышен на ту же сумму),
хотя ШАА и итоговая выручка клиники совпадают день в день. Проверили:
ни один врач не меняет \texttt{\seqsplit{department}} между 2 и 3 этажом в течение
года (не «переезд в середине года»), причина не найдена — возможно,
владелец вручную переносит выручку 1--2 врачей между этажами в своей
таблице по причине, не отражённой в ФОТ. Раз выручка используется как
делитель \texttt{\seqsplit{upr\_rate}}, эта неточность попадает и в распределение
Управления между 2 и 3 этажом.

\subsubsection{Маркетинг: 80\,\% психиатрия / 20\,\% психология (без изменений)}
\label{sec:overhead-mkt}

Не зависит от этажа — тот же принцип, что и в первой версии раздела.
Пул за месяц: ФОТ Маркетинга + его налоги (из \texttt{\seqsplit{realt\_payroll}} напрямую,
роль «ФОТ Маркетинг», без фильтра по врачу) + статьи «Рекламный бюджет
общий»/«Маркетинговые подрядчики»/«Телефония, связь, боты» из
\texttt{\seqsplit{realt\_pl\_by\_group\_month}} (группа «Маркетинг»). По решению
владельца 2026-09-19 пул делится 80\,\% на визиты психиатрии
(\texttt{\seqsplit{role\_group}} IN («ФОТ Психиатры», «ФОТ ШАА»)) и 20\,\% на визиты
психологии (\texttt{\seqsplit{role\_group}} = «ФОТ Психологи»); \textbf{психотерапевты
отдельно не выделяются} — своего \texttt{\seqsplit{role\_group}} у них нет (есть только
как \texttt{\seqsplit{doctor\_type}} в \texttt{\seqsplit{realt\_service\_categories}}, справочнике по
УСЛУГЕ, а не по сотруднику, нигде не используемом — заводить под это
вторую, непроверенную ось классификации не стали), входят в 20\,\%
психологии:

\[
  \texttt{\seqsplit{mkt\_psy\_rate}} = \frac{0.8 \times \text{Пул}_{\text{маркетинг}}}{\texttt{\seqsplit{visits\_psy}}}
  \qquad
  \texttt{\seqsplit{mkt\_pso\_rate}} = \frac{0.2 \times \text{Пул}_{\text{маркетинг}}}{\texttt{\seqsplit{visits\_pso}}}
\]

Каждому визиту присваивается ставка \textbf{его} \texttt{\seqsplit{role\_group}} — визит
психиатра получает \texttt{\seqsplit{mkt\_psy\_rate}}, психолога ---
\texttt{\seqsplit{mkt\_pso\_rate}}. Для «Все врачи» (без фильтра) это автоматически
даёт средневзвешенную по факту визитов ставку — специальной ветки для
этого случая в SQL нет, она получается сама из механизма атрибуции
(раздел~\ref{sec:overhead}).

\subsubsection{Прибыль после всех затрат}
\label{sec:overhead-profit}

\[
  \text{Прибыль после всех затрат} = \text{revenue} - \text{fot\_own} + \text{fot\_own\_taxes}
    + \texttt{\seqsplit{fixed\_attributed}} + \texttt{\seqsplit{mkt\_attributed}}
\]

Расширяет «Прибыль после ФОТ» (раздел~\ref{sec:monthly}) двумя новыми
слагаемыми (\texttt{\seqsplit{fixed\_attributed}}/\allowbreak\texttt{\seqsplit{mkt\_attributed}} — уже
просуммированная по видимым визитам атрибуция, отрицательный знак,
поэтому прибавляются, а не вычитаются — та же логика знаков, что и у
\texttt{\seqsplit{fot\_taxes}}).

\textbf{\texttt{\seqsplit{fot\_own}}/\texttt{\seqsplit{fot\_own\_taxes}} — это НЕ \texttt{\seqsplit{fot\_total}}/\allowbreak
\texttt{\seqsplit{taxes\_total}}} (существующая строка 22 карточки~183) — это ОТДЕЛЬНАЯ
пара колонок (сумма ФОТ/налогов ТОЛЬКО ролей Психиатры+Психологи+ШАА,
\textbf{всегда} без Админов/Управления/Маркетинга). Найдено и починено
2026-09-19 в тот же вечер: если использовать здесь \texttt{\seqsplit{fot\_total}}
напрямую, при \texttt{\seqsplit{doctor\_type}}=«Все» происходит \textbf{двойной счёт} —
\texttt{\seqsplit{fot\_total}} при «Все» и так включает ФОТ Админов+Управления+
Маркетинга (см. ГОЧТЯ раздела~\ref{sec:overhead-doctor-type-bug}), а
\texttt{\seqsplit{fixed\_attributed}} их \textbf{тоже} включает — на проде это занижало
«Прибыль после ФОТ и накладных» ровно на величину ФОТ Админов+Управления
(проверено: строка была на $\approx$2\,750~₽/визит ниже, чем должна).

\textbf{Налоги на накладные/маркетинг — не распределены отдельной долей от
маржи}, отложено по решению владельца 2026-09-19 (уже включённые
ФОТ-налоги Админов/Управления/Маркетинга — это НДФЛ+взносы с ИХ ФОТ, не
налог на прибыль всей клиники — она вообще исключена из накладных, см.
раздел~\ref{sec:overhead-upr}).

\textbf{Проверено на реальных данных (2026-09-19, финальная per-этаж
версия, запрос через Metabase API, год 2026, «Все врачи», ШАА
включена):} «Прибыль после ФОТ и накладных на визит» — Янв~681,
Фев~580, Мар~1\,549, Апр~548, Май~1\,183, Июн~732, Июл~$-433$, Авг~$-113$
— положительна почти весь год, минус только в июле-августе (небольшой).
Более ранняя (плоская, до этажей) версия давала стабильно отрицательные
$-2\,800$ до $-5\,500$~руб./визит — разница объясняется совокупно
контаминацией АПДШ/Шмилович, дублированием Р/с↔Начисления и переходом
на per-этаж методику (см. историю в начале раздела~\ref{sec:overhead}).
«Налог на прибыль/УСН» в этих цифрах \textbf{не учтён} (временно
исключён, см. выше) — реальная прибыль после налога ниже.

\textbf{ГОЧТЯ (проверено эмпирически 2026-09-19):} после добавления
строк 23--26 карточка~183 перестала выполняться в ClickHouse —
\code{Too many optimizations applied to query plan. Current limit 10000}
(\code{TOO\_MANY\_QUERY\_PLAN\_OPTIMIZATIONS}). 26-строчный \code{UNION ALL},
где КАЖДАЯ ветка обращается к CTE \texttt{\seqsplit{monthly}} (а через неё — к цепочке
из 6 новых CTE плюс существующие \texttt{\seqsplit{base}}/\texttt{\seqsplit{filtered}} с оконной
функцией), уже упирается в дефолтный лимит планировщика ClickHouse — до
добавления этих строк (24 ветки) лимита хватало. \textbf{Фикс} — явный
\code{SETTINGS query_plan_max_optimizations_to_apply = 100000} в конце
запроса (после \code{ORDER BY}), проверено на проде. Если будете
добавлять в эту карточку ещё строки — вероятно, лимит придётся поднимать
ещё, либо переписывать 26-way \code{UNION ALL} на другую конструкцию (не
делали — не тривиально, каждая ветка пивота требует своей агрегации по
12 месяцам).

\subsubsection{Побочный баг, найденный и починенный тем же вечером
2026-09-19: строки 12--17 пустели при фильтре «Тип врача»}
\label{sec:overhead-doctor-type-bug}

Независимо от раздела~\ref{sec:overhead} выше, при разборе с владельцем
выяснилось: строки «ФОТ Администраторы/Управление/Маркетинг ---
Проценты/Оклад+Бонус на визит» (12--17, СУЩЕСТВОВАЛИ ДО этой сессии, не
новые) были \textbf{полностью пустые} при любом значении параметра «Тип
врача» кроме «Все». Причина --- в \texttt{\seqsplit{payroll\_f}}
(раздел~\ref{sec:doctor-econ}) условие фильтра по \texttt{\seqsplit{doctor\_type}}
пропускало только \code{role_group IN ('ФОТ Психиатры', 'ФОТ ШАА')} или
\code{= 'ФОТ Психологи'} --- Администраторы/Управление/Маркетинг не
попадают НИ В ОДНУ из веток, поэтому выпадали из выборки целиком.
\textbf{Исправлено} --- добавлена третья ветка \code{OR role_group IN
('ФОТ Администраторы', 'ФОТ Управление', 'ФОТ Маркетинг')}: эти три роли
теперь \textbf{не подчиняются} фильтру «Тип врача» вообще (симметрично
разделу~\ref{sec:overhead}, где то же самое сделано для накладных) ---
доктором в смысле этого фильтра они не являются.

\textbf{Побочный эффект, тоже найден и починен в тот же заход:} \code{fot_total}/
\code{taxes_total} (используются в «ФОТ всего»/«Доля ФОТ в выручке»/«Прибыль
после ФОТ», раздел~\ref{sec:monthly}) раньше складывали ВСЕ строки
\texttt{\seqsplit{payroll\_f}} --- после починки строк 12--17 это означало бы, что
полный клиники-целиком накладной ФОТ (посчитанный на весь объём визитов
клиники) начинает делиться на визиты ОДНОГО типа врача (знаменатель
резко меньше). Проверено на проде: «Прибыль после ФОТ» на визит для
«Психологические» за год уходила в $-8\,627$~руб. --- искусственно, не
отражая реальность. \textbf{Исправлено вторым шагом}: \code{fot\_total\_pct}/\allowbreak
\code{fot\_total\_oklad}/\code{taxes\_total} в \texttt{\seqsplit{payroll\_pivot}} теперь ЯВНО
вычитают Админов/Управление/Маркетинг, когда \texttt{\seqsplit{doctor\_type}} \textbf{не}
«Все» (при «Все» вычитать нечего --- вычитание равно 0, поведение не
меняется). Итог: строки 12--17 показывают полный клиники-целиком
накладной ФОТ при любом фильтре (как просил владелец), а «ФОТ
всего»/«Прибыль после ФОТ» остаются посчитаны ТОЛЬКО по клиническому ФОТ
выбранного типа врача (как было до этой правки, семантика сохранена).
Проверено на проде повторно после фикса: «Прибыль после ФОТ» для
«Психологические» вернулась к разумным значениям ($\approx 5\,879$ /
$7\,133$~руб.\ янв./год).

\section{P\&L по группам (\texttt{realt\_pl\_by\_group\_month})}
\label{sec:pl}

Объединение (\code{UNION ALL}, не \texttt{\seqsplit{JOIN}}) трёх источников ---
\texttt{\seqsplit{realt\_bank\_account}}, \texttt{\seqsplit{realt\_cash}}, \texttt{\seqsplit{realt\_accruals}} --- по
общей роли колонок (\texttt{\seqsplit{pl\_article}}, дата операции/начисления, сумма),
плюс перенесённые из \texttt{\seqsplit{realt\_payroll}} строки ФОТ Маркетинга и
Управления (переезжают в свою группу, а не остаются отдельной строкой
«ФОТ»). Whitelist --- 24 статьи, единый источник для обеих Metabase
Model, которые раньше дублировали этот список (Model~178 и карточка~179).

\begin{longtable}{@{}L{3.6cm}L{11cm}@{}}
\caption{Группировка статей в \texttt{realt\_pl\_by\_group\_month}}
\label{tab:pl-groups}\\
\toprule
Группа & Статьи \\
\midrule
\endfirsthead
\toprule
Группа & Статьи \\
\midrule
\endhead
\bottomrule
\endfoot
\textbf{Помещение} & Аренда, Ком.\ услуги, Санпэдрежим/охрана труда --- каждая
  на 2 этаже, 3 этаже и в ШАА (9 статей) \\
\textbf{Маркетинг} & Телефония/связь/боты, Рекламный бюджет общий,
  Маркетинговые подрядчики (3 статьи из источников) \newline
  + \code{ФОТ Маркетинг} = $-\sum \texttt{\seqsplit{accrued\_total}}$ по \code{role='ФОТ Маркетинг'}
  \newline
  + \code{Взносы и НДФЛ ФОТ Маркетинг} = $\sum \texttt{\seqsplit{ndfl}} + \sum \texttt{\seqsplit{contributions}}$
  той же роли \\
\textbf{Административные} & \code{ФОТ Управление} = $-\sum \texttt{\seqsplit{accrued\_total}}$
  по \code{role='ФОТ Управление'} \newline
  + \code{НДФЛ - Управление} и \code{Взносы ФОТ - Управление} той же роли
  (отдельными строками, не суммой) \newline
  + 8 статей из источников (Сервисы и подписки, Интернет/телефония/почта,
  Ремонт/Оборудование/Мебель, Канцелярия, Аутсорс/консалтинг/юристы,
  Мероприятия/обучение/подарки, Банковские услуги, Прочие админ расходы) \\
\textbf{Внегрупповые} & Прочие доходы, Проценты по кредитам, Налог на
  прибыль / УСН, Амортизация --- каждая сама себе «группа», не
  суммируется ни с чем \\
\end{longtable}

\textbf{Знаки:} суммы из источников (\texttt{\seqsplit{amount}}/\allowbreak\texttt{\seqsplit{accrual\_amount}})
--- как в исходном регистре (обычно отрицательные для расходов). Строки
ФОТ намеренно \textbf{негируются}
($-\sum \texttt{\seqsplit{accrued\_total}}$), чтобы попасть в P\&L с тем же знаком,
что и остальные статьи-расходы; \texttt{\seqsplit{ndfl}}/\allowbreak\texttt{\seqsplit{contributions}} уже
отрицательные в источнике, не негируются повторно.

\subsection{Два фикса 2026-09-19 (найдены при построчной сверке с P\&L владельца)}
\label{sec:pl-fixes-2026-09-19}

\textbf{1. АПДШ/«Шмилович личное» исключены всегда.} У
\texttt{\seqsplit{realt\_bank\_account}}/\texttt{\seqsplit{realt\_cash}} есть колонка \texttt{\seqsplit{project}}, у
\texttt{\seqsplit{realt\_accruals}} — \texttt{\seqsplit{tag}} (три разных поля, один смысл: этаж/
направление — «2 этаж»/«3 этаж»/«ШАА»/«АПДШ»/«Шмилович личное»). VIEW
эту колонку не читала вообще — «АПДШ» (отдельное направление/юрлицо, не
клиника Реальт) и «Шмилович личное» (личные расходы, не бизнес) утекали
в статьи whitelist наравне с клиникой: на реальных данных найдено
$\approx$2.9\,млн~₽ АПДШ + $\approx$21\,тыс~₽ Шмилович личное за 2026 год
только в «Банковские услуги»/«Аутсорс, консалтинг, юристы»/«Налог на
прибыль / УСН». \textbf{Исправлено} — оба тега исключены из
\texttt{\seqsplit{expense\_combined}} всегда, независимо ни от каких тумблеров.

\textbf{2. \texttt{\seqsplit{realt\_bank\_account}} участвует в P\&L только строками с
заполненной \texttt{\seqsplit{accrual\_date}}.} Раньше формула была
\code{coalesce(accrual_date, operation_date)} — при отсутствующей
\texttt{\seqsplit{accrual\_date}} падала на дату операции. Оказалось: у части операций
(в основном многомесячные лицензии/подписки) в Р/с нет
\texttt{\seqsplit{accrual\_date}} \textbf{потому что} эта же операция уже размазана по
месяцам во вкладке «Начисления» отдельными строками — «Начисления» и
есть механизм разбивки денежной операции по периодам/суммам/статьям
(объяснение владельца, 2026-09-19). Без фильтра сумма считалась
\textbf{дважды}: целиком в Р/с на дату платежа, и ещё раз размазанной по
Начислениям. Конкретные найденные пары (сумма Р/с делится ровно на N
Начисления-строк по N месяцам): Понамарёв/«ПрофитПросто» (80\,000~₽ →
6$\times$13\,333), Ваззап (32\,400 → 6$\times$5\,400), Яндекс~360, СКБ
Контур, Клиентикс~ERP, 1С:Фреш, АДВЕРТМЕД/Битрикс, а также банковские
комиссии «Точка Банк» и часть налоговых/страховых платежей.
\textbf{Исправлено} — \texttt{\seqsplit{realt\_bank\_account}} берёт строго
\texttt{\seqsplit{accrual\_date}} (без fallback), строки без неё выпадают из этого
источника целиком (доверяем, что их отразят Начисления).

\textbf{\texttt{\seqsplit{realt\_cash}} сознательно НЕ тронут} тем же фильтром —
проверено эмпирически: у него \texttt{\seqsplit{accrual\_date}} не заполнена НИ РАЗУ (0
из 256 строк за всю историю), задвоения с Начислениями для него не
нашли (мелкие валютные списания типа Pixelcut/Gamma.app/Combot/Freepik —
легитимные разовые траты без пары в Начислениях). Попытка применить
фильтр и к Наличным ухудшила сверку (март «Телефония, связь, боты» ушёл
с верных 73\,807~₽ на заниженные 61\,872~₽) — оставлен на
\code{coalesce(accrual\_date, operation\_date)}, что для этой таблицы
эквивалентно всегда \texttt{\seqsplit{operation\_date}}.

\textbf{Результат сверки после обоих фиксов} (112 ячеек: 14 статей
групп «Маркетинг»+«Административные» $\times$ 8 месяцев, сверено против
таблицы владельца построчно) — 104 совпали день в день, 8 разошлись на
величину, зеркальную между соседними месяцами (одна и та же сумма
«утекла» в другой месяц у одной из сторон — например, «Рекламный
бюджет общий» июль/август $-14\,500$/$+14\,500$) — мелкая настройка
границ месяца, не системная ошибка.

\textbf{\texttt{\seqsplit{amount\_ex\_shaa}}} (новая колонка, единообразно с
\texttt{\seqsplit{realt\_expenses\_by\_month}}) --- сумма без строк с
\texttt{\seqsplit{project}}/\texttt{\seqsplit{tag}}=«ШАА» (\texttt{\seqsplit{is\_shaa}} вычисляется из этого же
поля). Заменила текстовый поиск «ШАА» в названии статьи, который раньше
использовался только для группы «Помещение» — для
Административных/Маркетинга/Внегрупповых тумблер «Учитывать ШАА» вообще
не работал (найдено: тег ШАА на «Налог на прибыль / УСН»
$\approx$914\,тыс~₽/год учитывался независимо от тумблера).

\section{Прочие витрины}
\label{sec:other-views}

\subsection{\texttt{realt\_expenses\_by\_month}}

\[
  \texttt{\seqsplit{amount}} = \sum \texttt{\seqsplit{realt\_expenses.amount}}
\]
\[
  \texttt{\seqsplit{amount\_ex\_shaa}} = \sum \texttt{\seqsplit{realt\_expenses.amount}} \text{ где } \texttt{\seqsplit{is\_shaa}} = 0
\]

Группировка --- месяц~$\times$~\texttt{\seqsplit{expense\_type}} (набор конкретных
статей 2025~$\ne$~2026, поэтому группируем по типу, не по статье).
\texttt{\seqsplit{amount\_ex\_shaa}} --- для среза «без Шмиловича», согласованного по
смыслу с \texttt{\seqsplit{revenue}} (раздел~\ref{sec:monthly}), хоть и не тем же
самым методом exclusion (\texttt{\seqsplit{is\_shaa}} --- признак статьи расхода,
проставленный вручную в источнике, а не текстовое/личностное правило).

\subsection{\texttt{realt\_revenue\_by\_service}}

\[
  \texttt{\seqsplit{revenue}} = \sum \texttt{\seqsplit{amount}}
  \qquad
  \texttt{\seqsplit{visits}} = \text{count}(*)
  \qquad
  \texttt{\seqsplit{avg\_check}} = \frac{\texttt{\seqsplit{revenue}}}{\texttt{\seqsplit{visits}}}
\]

Группировка --- месяц~$\times$~\texttt{\seqsplit{service}}. Фильтры --- те же базовые
(раздел~\ref{sec:base-filters}), включая текстовое исключение
Шмиловича/Онегиной/«тест».

\section{Параметры (фильтры) дашбордов}
\label{sec:filters}

Общие для дашборда «Помесячная юнитка (по врачам, корректная)» (id~7,
карточки 183/184) параметры --- не формулы метрик, а то, что сужает
выборку визитов/зарплат \textbf{до} расчёта метрик из
разделов~\ref{sec:doctor-econ}.

\begin{longtable}{@{}L{3.3cm}L{4.2cm}L{6.1cm}@{}}
\caption{Параметры дашборда id~7}
\label{tab:dash-filters}\\
\toprule
Параметр & Значения & Логика \\
\midrule
\endfirsthead
\toprule
Параметр & Значения & Логика \\
\midrule
\endhead
\bottomrule
\endfoot
Год & число & \code{toYear(month) = год} \\
Тип клиента & Все / Первичный / Повторный &
  По \texttt{\seqsplit{mvisits}} = число визитов ЭТОГО \texttt{\seqsplit{card\_number}} \textbf{в
  этом же месяце} (\code{count(*) OVER (PARTITION BY card_number, month)}).
  Первичный --- \code{mvisits = 1}, Повторный --- \code{mvisits >= 2}.
  \textbf{Не путать} с «Новым клиентом» (раздел~\ref{sec:terms}) --- там
  речь о первом визите за всю историю, здесь --- про визиты внутри
  месяца. \\
Тип врача & Все / Психиатрические / Психологические &
  Психиатрические = \code{role_group IN ('ФОТ Психиатры', 'ФОТ ШАА')};
  Психологические = \code{role_group = 'ФОТ Психологи'}. \\
Учитывать ШАА & Да / Нет &
  Нет $\Rightarrow$ строки с \code{role_group = 'ФОТ ШАА'} исключаются
  целиком из выборки (и из визитов, и из зарплаты). \\
Показатели на & Визит / Клиент &
  Переключает знаменатель у всех денежных строк дашборда между
  \texttt{\seqsplit{visits}} и \texttt{\seqsplit{clients}} (см.\ раздел~\ref{sec:doctor-econ}). \\
Врач & дропдаун (82 значения) &
  Сужает выборку до одного \texttt{\seqsplit{employee\_id}} (через сопоставление по
  \texttt{\seqsplit{doctor\_name}}), значения --- из отдельной карточки-справочника
  (\texttt{\seqsplit{realt\_doctors\_list.sql}}, id~185). \\
\end{longtable}

\section{Схема соединения таблиц}
\label{sec:er}

\begin{figure}[p]
\centering
\resizebox{\linewidth}{!}{%
\begin{tikzpicture}[
  node distance=14mm and 16mm,
  every node/.style={font=\small},
  tablebox/.style={draw, rounded corners, align=left, inner sep=6pt, fill=gray!4},
  tablebox2/.style={tablebox, draw=gray, dashed, fill=gray!2, text=gray!60!black},
  arr/.style={-{Latex[length=2.5mm]}, thick},
  arr2/.style={-{Latex[length=2.5mm]}, thick, dashed, gray}
]

\node[tablebox] (kl) {%
  \begin{tabular}{@{}l@{}}
    \textbf{klientiks\_operations}\\
    \hline
    card\_number\\
    doctor\\
    service\\
    visit\_start\\
    amount
  \end{tabular}%
};

\node[tablebox, right=26mm of kl] (emp) {%
  \begin{tabular}{@{}l@{}}
    \textbf{realt\_employees}\\
    \hline
    employee\_id\\
    full\_name
  \end{tabular}%
};

\node[tablebox, right=of emp] (pay) {%
  \begin{tabular}{@{}l@{}}
    \textbf{realt\_payroll}\\
    \hline
    employee\_id\\
    period\\
    role, department\\
    pay\_type\\
    accrued\_total\\
    ndfl, contributions
  \end{tabular}%
};

\draw[arr] (kl) -- (emp) node[midway, above, font=\scriptsize, text width=24mm, align=center] {doctor $\approx$ full\_name\\ (регистронезависимо)};
\draw[arr] (emp) -- (pay) node[midway, above, font=\scriptsize] {employee\_id};

\node[tablebox2, below=of kl] (svc) {%
  \begin{tabular}{@{}l@{}}
    \textbf{realt\_service\_categories}\\
    \hline
    service\\
    doctor\_type\\
    qualification
  \end{tabular}%
};
\draw[arr2] (kl) -- (svc) node[midway, right, font=\scriptsize, text width=3.6cm] {service (ключ есть, но не используется ни в одной VIEW)};

\node[tablebox, below=of svc, xshift=42mm, yshift=-8mm] (exp) {%
  \begin{tabular}{@{}l@{}}
    \textbf{realt\_expenses}\\
    \hline
    period, article\\
    expense\_type\\
    amount, is\_shaa
  \end{tabular}%
};
\node[tablebox, right=of exp] (bank) {%
  \begin{tabular}{@{}l@{}}
    \textbf{realt\_bank\_account}\\
    \hline
    pl\_article\\
    accrual\_date /\\
    \ operation\_date\\
    amount, accrual\_amount
  \end{tabular}%
};
\node[tablebox, right=of bank] (cash) {%
  \begin{tabular}{@{}l@{}}
    \textbf{realt\_cash}\\
    \hline
    pl\_article\\
    accrual\_date /\\
    \ operation\_date\\
    amount, accrual\_amount
  \end{tabular}%
};
\node[tablebox, right=of cash] (acc) {%
  \begin{tabular}{@{}l@{}}
    \textbf{realt\_accruals}\\
    \hline
    pl\_article\\
    accrual\_date\\
    amount, accrual\_amount
  \end{tabular}%
};

\node[draw=gray, dashed, rounded corners, fit=(exp)(bank)(cash)(acc), inner sep=7mm,
      label={[gray, font=\scriptsize]above:Источники P\&L --- НЕ джойнятся друг с другом, объединяются UNION ALL по pl\_article (realt\_pl\_by\_group\_month)}] {};

\end{tikzpicture}%
}
\caption{Таблицы-источники Реальта и ключи их соединения. Пунктир ---
связь, которая существует по данным, но пока не задействована ни в
одной VIEW/карточке (раздел~\ref{sec:sources}).}
\label{fig:realt-er}
\end{figure}

\textbf{Как читать:} сплошная стрелка --- реальный \texttt{\seqsplit{JOIN}} внутри
VIEW (раздел~\ref{sec:categorized}), подпись --- условие соединения.
Пунктирная стрелка --- потенциальный ключ (есть в обеих таблицах), но
ни одна текущая VIEW/карточка его не использует. Пунктирная рамка внизу
--- три источника P\&L, которые не связаны JOIN'ом друг с другом вообще:
каждая строка одного источника не имеет пары в двух других, они просто
складываются построчно (раздел~\ref{sec:pl}).

Что из чего строится (VIEW $\to$ таблицы-источники) --- сводно:

\begin{itemize}[leftmargin=*]
\item \texttt{\seqsplit{realt\_metrics\_by\_month}} (раздел~\ref{sec:monthly}) ---
  \textbf{с 2026-09-19}: \texttt{\seqsplit{realt\_visits\_categorized}} (визиты, категоризация
  по \texttt{\seqsplit{role\_group}}, личностно) + \texttt{\seqsplit{realt\_payroll}} (ФОТ); раньше ---
  \texttt{\seqsplit{klientiks\_operations}} напрямую + \texttt{\seqsplit{realt\_payroll}}, категоризация
  по тексту услуги (см.\ раздел~\ref{sec:base-filters});
\item \texttt{\seqsplit{realt\_revenue\_by\_service}} --- \texttt{\seqsplit{klientiks\_operations}}
  (\textbf{всё ещё} по тексту услуги, не переведена на
  \texttt{\seqsplit{realt\_visits\_categorized}} --- см.\ раздел~\ref{sec:other-views});
\item \texttt{\seqsplit{realt\_expenses\_by\_month}} --- \texttt{\seqsplit{realt\_expenses}};
\item \texttt{\seqsplit{realt\_pl\_by\_group\_month}} (раздел~\ref{sec:pl}) ---
  \texttt{\seqsplit{realt\_bank\_account}} + \texttt{\seqsplit{realt\_cash}} + \texttt{\seqsplit{realt\_accruals}} +
  \texttt{\seqsplit{realt\_payroll}} (только роли Маркетинг/Управление);
\item \texttt{\seqsplit{realt\_visits\_categorized}} --- \texttt{\seqsplit{klientiks\_operations}} +
  \texttt{\seqsplit{realt\_employees}} + \texttt{\seqsplit{realt\_payroll}} (только для определения
  \texttt{\seqsplit{role\_group}}, ПО МЕСЯЦУ визита --- см.\ раздел~\ref{sec:categorized} ---
  сами суммы зарплаты не читает);
\item \texttt{\seqsplit{realt\_payroll\_categorized}} --- только \texttt{\seqsplit{realt\_payroll}} (с
  2026-09-19 без доп.\ подзапроса --- \texttt{\seqsplit{department}} уже есть в своей же
  строке);
\item \texttt{\seqsplit{realt\_doctor\_month}} / \texttt{\seqsplit{realt\_role\_month}}
  (раздел~\ref{sec:doctor-econ}) --- \texttt{\seqsplit{realt\_visits\_categorized}} +
  \texttt{\seqsplit{realt\_payroll\_categorized}} + \texttt{\seqsplit{realt\_employees}} (для ФИО).
\end{itemize}

\textit{Диаграмма выше (рис.~\ref{fig:realt-er}) отражает состояние на
момент её отрисовки (сплошная стрелка klientiks\_operations
$\to$ realt\_employees $\to$ realt\_payroll) --- сам факт JOIN'а не
изменился 2026-09-19, изменилось только условие (по месяцу, а не
\texttt{\seqsplit{employee\_id}} без учёта периода) --- перерисовывать не требуется.}

\section{Где что лежит}

\small
\begin{longtable}{@{}L{4cm}L{7.3cm}L{4cm}@{}}
\caption{Соответствие: раздел документа $\to$ файл в репозитории $\to$
Metabase}
\label{tab:where}\\
\toprule
Раздел & Файл в репозитории & Metabase \\
\midrule
\endfirsthead
\toprule
Раздел & Файл в репозитории & Metabase \\
\midrule
\endhead
\bottomrule
\endfoot
\ref{sec:monthly}--\ref{sec:ratio-yearly} Месячные метрики,
  полная\,/\,ШАА\,/\,скользящая выручка, «За год» для ratio &
  \texttt{\seqsplit{schema\_realt\_metrics\_views.sql}}
  (\texttt{\seqsplit{realt\_metrics\_by\_month}}); \texttt{\seqsplit{realt\_metrics\_model.sql}};
  \texttt{\seqsplit{realt\_monthly\_metrics.sql}} (формула «За год» для ratio-строк ---
  только здесь) & Модель id~98, карточка id~186, дашборд «Ежемесячные
  метрики» id~8 (Таблица, осознанное исключение из правила Визуалов) \\
\ref{sec:categorized}--\ref{sec:doctor-econ} Юнит-экономика по врачу &
  \texttt{\seqsplit{schema\_realt\_metrics\_views.sql}} (\texttt{\seqsplit{realt\_visits\_categorized}},
  \texttt{\seqsplit{realt\_payroll\_categorized}}, \texttt{\seqsplit{realt\_doctor\_month}},
  \texttt{\seqsplit{realt\_role\_month}}); \texttt{\seqsplit{realt\_unitka\_by\_doctor.sql}};
  \texttt{\seqsplit{realt\_doctors\_table.sql}} & Карточки 183/184, дашборд id~7 \\
\ref{sec:overhead} Накладные и маркетинг на визит/клиента &
  \texttt{\seqsplit{realt\_unitka\_by\_doctor.sql}} (строки 23--26);
  \texttt{\seqsplit{realt\_doctors\_table.sql}} (колонки «Накладные на юнит»/«Маркетинг
  на юнит»/«Прибыль после всех затрат») — формулы только в native SQL
  карточек, отдельной VIEW нет & Карточки 183/184, дашборд id~7 \\
\ref{sec:pl} P\&L по группам & \texttt{\seqsplit{schema\_realt\_metrics\_views.sql}}
  (\texttt{\seqsplit{realt\_pl\_by\_group\_month}}) & Модель id~178, карточка id~179 \\
\ref{sec:other-views} Расходы / выручка по услугам &
  \texttt{\seqsplit{schema\_realt\_metrics\_views.sql}}
  (\texttt{\seqsplit{realt\_expenses\_by\_month}}, \texttt{\seqsplit{realt\_revenue\_by\_service}}) & --- \\
\end{longtable}
\normalsize

\section{Как поддерживать документ в актуальном состоянии}
\label{sec:maintain}

\begin{itemize}[leftmargin=*]
\item Формула метрики изменилась в \texttt{\seqsplit{schema\_realt\_metrics\_views.sql}}
  или в одной из карточек \texttt{\seqsplit{src/ metabase\_queries/ realt\_*.sql}} ---
  поправить соответствующий раздел этого файла \textbf{тем же коммитом}.
\item Появилась новая метрика/VIEW --- добавить раздел по образцу уже
  существующих (формула + таблица частных случаев, если их много) и
  строку в таблицу~\ref{tab:where}.
\item \textbf{Новая метрика --- ratio (среднее/доля/X на клиента), не
  сумма}: (1) явно написать в этом файле, что колонка --- ratio и как
  считается её корректный годовой/итоговый агрегат (по образцу
  раздела~\ref{sec:ratio-yearly}); (2) `COMMENT COLUMN` в
  \texttt{\seqsplit{schema\_realt\_metrics\_views.sql}} должен предупреждать о том же для
  бота; (3) если карточка-Таблица с такой строкой ложится на дашборд
  напрямую (не как Визуал) --- это осознанное исключение из правила
  Metabase-стандартов, отметить в \texttt{\seqsplit{docs/architecture-map.md}} (раздел
  про Metabase) и в самих `architecture-standarts.md`, если такого
  прецедента там ещё нет.
\item \texttt{\seqsplit{realt\_employees}}/\allowbreak\texttt{\seqsplit{realt\_service\_categories}}
  (раздел~\ref{sec:sources}) --- на момент написания документа код их
  создания не смёржен в \texttt{\seqsplit{master}} (ветки \texttt{\seqsplit{feat/realt-employees}},
  \texttt{\seqsplit{feat/realt-service-categories}}), хотя таблицы задеплоены и
  используются в проде. Смёржить ветки --- отдельная задача, не блокирует
  актуальность формул здесь (они описывают то, что реально в проде).
\end{itemize}

\end{document}
